% Isolated exploratory note.  Not part of the frozen CP31 paper.
\subsection*{Two elementary tests for a reduced Pfaffian incidence form}

Let \(k\) be algebraically closed of characteristic two, let \(N\ge3\)
be odd, and let \(B(\lambda)\) be an \(N\)-by-\(N\) alternating
matrix of linear forms in \(\lambda_0,\ldots,\lambda_{N-1}\).
Assume its generic rank is \(N-1\), and divide its maximal-Pfaffian
vector by its common polynomial factor.  Denote the resulting
coprime vector by \(\sigma(\lambda)\), and assume its projective
map is birational.  Put
\[
                    q(\lambda)=\sum_i\lambda_i\sigma_i(\lambda)^2.
\]
The following direct observations are candidates awaiting independent
audit.  They make no assertion about the base codimension for the
actual Cartier pencil, nor about the common Pfaffian factor.

\paragraph{No linear factor.}
The polynomial \(q\) has no linear factor.  Indeed, make a constant
linear change in \(\lambda\) and the corresponding square-root
contragredient change in the radical coordinates so that the
putative factor is \(\lambda_0\).  Such a change exists over the
perfect field and preserves the linear alternating pencil and
the displayed pairing.  Setting \(\lambda_0=0\), the distinct
parity classes of the monomials in
\(\sum_{i>0}\lambda_i\sigma_i^2\) show that every
\(\sigma_i\), \(i>0\), is divisible by \(\lambda_0\).
Write \(\sigma_i=\lambda_0\tau_i\).  Coprimality implies that
\(\sigma_0\) is not divisible by \(\lambda_0\).

The identity \(B\sigma=0\) then forces
\(B_{i0}=a_i\lambda_0\) with constants \(a_i\), since
\(B_{i0}\) is linear.  Its zeroth row gives
\(\sum_{i>0}a_i\tau_i=0\), hence the constant linear relation
\(\sum_{i>0}a_i\sigma_i=0\).  Birationality makes the coordinate
functions \(\sigma_i\) linearly independent, so all \(a_i=0\).
The zeroth row and column of \(B\) vanish identically.  Its unique
generic radical is then the constant zeroth coordinate line,
contradicting birationality.  This proves the observation.

\paragraph{A conditional irreducibility test.}
If the common zero locus of the reduced coordinates \(\sigma_i\)
has codimension at least three in projective space, then \(q\)
is irreducible.  In fact
\(\partial q/\partial\lambda_i=\sigma_i^2\), so its singular
support is exactly that common zero locus.  Coprimality rules out
any repeated irreducible factor of \(q\).  If two distinct
components existed, their intersection would have codimension two
and lie in the singular support, a contradiction.  The codimension
hypothesis is not inferred merely from generic maximal rank or
from birationality of \(\sigma\).

\paragraph{Remaining distinction.}
Even this conditional irreducibility statement does not prove that
\(q=0\) meets the open where the original alternating pencil has
rank \(N-1\): its removed common Pfaffian factor could still vanish
on that hypersurface.  Nor does it exclude extra components of the
Lang incidence over lower-rank strata.  These are precisely the
additional geometric conditions that must be checked before using
the local wall theorem or a connectedness argument.

\paragraph{Primary-source comparison, not an imported theorem.}
The publisher original of Russo and Simis, \emph{On Birational Maps
and Jacobian Matrices}, DOI \(10.1023/A:1017572213947\), was opened
and its full Proposition 2.2 proof and Example 2.5 argument on
printed pages 342--344 were read.  Example 2.5 explicitly assumes
height three and the stated local generator/Fitting-ideal conditions;
it does not infer those hypotheses from birationality.  Those
conditions have not been verified for the present higher-degree
Cartier pencil.  No result from that example is imported here.
The adjacent monomial and Jacobian results have characteristic-zero
qualifications where printed, which are not suppressed.
No numerical computation was used in these two elementary arguments.
